\largerpage
\scp{“Eastern Indonesia” as a complex ethnographic transition zone}{02-06}
Following on from the discussion of contact-induced change in
the Austronesian languages of the region (\srf{02-03}),
it is important to note that evidence of contact is not limited
to language, but is also found in a range of distinct socio-cultural features
that have been observed for nearly 200 years \citep{Fox2015}.
\citet[1]{vanWouden1935} defined Eastern Indonesia as a distinct
ethnographic region “based upon various points of similarity
from an ethnographical point of view”.
The similarities he noted included a clan system,
exclusive cross-cousin marriage, and sociocosmic dualism,
collectively being features that set these societies apart from
Austronesian societies to the west.
\citet{Fox1980,Fox1989,Fox2015} revisited van Wouden’s work,
expanding the list of distinguishing features to include:

\begin{itemize}
	\item Concern with origins;
	\item Use of dual categories of complementary opposition;
	\item Identification of life processes and relationships in botanical metaphors; (also \citealp{Fox1971});
	\item The notion of ‘person’ as composed of \it{blood / flesh} and \it{semen / bone};
	\item Diarchy predicated on delegation or usurpation of authority;
	\item “House” as a primary descent group;
	\item Clans form political or ceremonial groupings;
	\item Marriage alliances link houses/clans;
	\item Alliances are viewed as the “flow of life”;
	\item Extended series of life-enhancing rituals joining predecessor to descendant in a cyclical translation of life.
\end{itemize}

Looking eastward into Melanesia
and south into northern Australia,
\citet[18]{Grimes2010} concludes “Eastern Indonesia is a zone where
social interaction has resulted in shared cultural heritages
with other Austronesian societies as well as non-Austronesian
societies.” Her conclusion for \ili{Buru} society (in central Maluku)
is that it could be described as an Austronesian speaking society
with a significant number of features common to societies speaking
Papuan languages in the greater Melanesian region,
including in the highlands of Papua New Guinea
and islands \ili{east} of New Guinea.
Whereas \ili{Fataluku} (in Timor-Leste) is “a society speaking a Papuan
language with many cultural features in common with their Austronesian
speaking neighbors \citep{McWilliam2007}.” From an ethnographic perspective,
there are additional examples supporting the view that eastern
Indonesia should be considered as a complex zone of transition \citep[191]{Fox2015}.
While some of these features by themselves are found in many
parts of the world,
this unique bundling of cultural features is not widespread in
Austronesian societies west of the Blust line.

More specifically, \citet{Fox2015} finds that the distribution
in terminologies and semantics relating to opposite sex siblings
and to the term for `child’s spouse’s parent' align roughly with the Blust line.
These are touched on in more detail in subsequent chapters.

While *mamah ‘father’s brother’ is reconstructed to \tsc{PAn},
in Wallacea many languages have developed a lexicalised term
for ‘mother’s brother’, not found in the west.
In some Wallacean languages the term is cognate with
\citeauthor{Stresemann1927}'s (\citeyear{Stresemann1927}) PAMS **məmə ‘mother’s brother’.
Other terms for ‘mother’s brother’ are also found.
This is an example of a \it{category innovation} (not a lexical
innovation with formal cognates),
that is likely to arise from substrate influence with variation in different regions.
This concept is actually much more nuanced.
For example, \ili{Buru} (central Maluku) \it{meme} ‘mother’s brother’
is a reflex of PAMS **məmə ‘mother’s brother’.
\ili{Tii} (Rote Island) \it{toʔo} ‘mother’s brother’ is not (it is
from *tau ‘person’ [\citealp[370]{Edwards2021}]),
but captures exactly the same concept or category in the kinship system.
Furthermore both languages use the same botanical metaphor from
‘tree, base of tree, source’ to point to the ‘mother’s brother with
the most senior rank, eldest’.
So \ili{Buru} \it{meme lahi {\tl} mem-lahin} ‘mother’s oldest (classificatory)
brother (lit: root/source uncle)’ (\it{lahin} ‘root; tree,
plant; base, source, cause’), is not formally cognate with,
but has almost identical semantics to,
\ili{Tii} \it{toʔo huu-k} ‘mother’s oldest (classificatory) brother
(lit: root/source uncle)’ (\it{huu-k} ‘root; tree,
plant; base, source, cause’; cf.
verbal doublet \it{ihuu-ioka} ‘take root’).
The category of ‘mother’s brother’ is widespread,
even though only some of the terms are formally cognate.
For example, \ili{Amarasi} \it{baba-f} ‘mother’s brother,
father-in-law’, \ili{Helong} \it{baki} ‘mother’s brother’,
\ili{Kemak} \it{naa-r} ‘mother’s brother (from *ñaRa ‘brother (f.s.))’,
\ili{Raklungu} \it{mimi} ‘mother’s brother’,
\ili{Selaru} \it{memi} ‘mother’s brother,
father-in-law’ reflect a very stable concept across a number of subgroups.

Further evidence is found in patterns of musical forms.
\citet{Blench2021} notes that polyphony is found throughout the
Austronesian cultures of Formosa (except one),
but was apparently left behind there when the Austronesian speaking
people moved out of Formosa.
However, vocal and instrumental polyphony is also widespread
where Papuan languages are spoken.
Referring to forms of polyphony found in Austronesian speaking
societies in eastern Flores
and parts of Timor,
Blench finds it “plausible that patterns found in eastern Indonesia
reflect the Papuan substrate”.

There is an abundance of direct evidence in regions where Papuan
languages are still spoken today,
and an abundance of trace evidence where there are no Papuan
languages spoken today, that allow us to posit contact as the source or trigger for a
number of features in the lexicons,
grammars, and even in the cultures of Austronesian-speaking peoples.
These features are found in the Papuan languages.
They are not generally found in the Austronesian languages west
of the Blust line.
The nature of the evidence implies intense contact,
bilingualism and linguistic and cultural accommodation over a
long period of time.
However they do not address the issues of whether or not there
was a shared parent language,
nor to subgrouping issues in the Austronesian languages in the region.
